\documentclass[10pt,a4paper]{amsart}
\usepackage[margin=19mm]{geometry}
\usepackage{amsmath,amssymb,mathtools}
\usepackage[scr=rsfs,cal=euler]{mathalfa}
\usepackage{xfrac}
\usepackage[unicode,hidelinks,bookmarks=false]{hyperref}
\newcommand{\ie}{i.e.\ }
\newcommand{\cf}{cf.\ }
\newcommand{\resp}{respectively\ }
\newcommand{\Subsection}{\subsection}
\providecommand{\nobreakdash}{\nobreak\hbox{-}}
\setlength{\emergencystretch}{2em}
\multlinegap0pt
\usepackage{amssymb}
\usepackage{mathtools}
\usepackage[shortlabels]{enumitem}
\newtheorem{theorem}{Theorem}
\newcommand{\RR}{\mathbb{R}}
\DeclareMathOperator{\col}{col}
\DeclareMathOperator{\rank}{rank}
\newcommand{\so}{\mathfrak{so}}
\title{Random Dot Product Graphs as Dynamical Systems: \\ Limitations and Opportunities}
\author{Giulio Valentino Dalla Riva}
\date{Source: arXiv:2603.05703v1 (CC BY 4.0)}
\begin{document}
\maketitle

\noindent\emph{Source and licence.} Random Dot Product Graphs as Dynamical Systems: \\ Limitations and Opportunities. Version arXiv:2603.05703v1 (CC BY 4.0), distributed by arXiv under the Creative Commons Attribution 4.0 International licence, http://creativecommons.org/licenses/by/4.0/, as recorded in the arXiv licence metadata for this exact version and shown on the abstract page https://arxiv.org/abs/2603.05703v1. Full licence notice: \texttt{licenses/arxiv-2603.05703-CC-BY-4.0.txt}.\par\smallskip

\noindent\textit{Editorial scope.} Counted unit: the article's theorem thm:invisible (source0.tex lines 289--291). The article's complete original proof of it is delivered with it (source0.tex lines 293--317, one \texttt{proof} environment of 25 lines). No mathematics is written, repaired or re-derived: the statement and the proof are the article's own lines, in the article's own notation, and no retained line is substituted or re-flowed.  No further source-local context is needed: every cross-reference of every retained line resolves inside the two delivered spans.  Nothing was omitted from any retained span, so the package carries no omission record.  No retained line contains a citation, so the wrapper carries no bibliography and no bibliography entry.  No retained line contains a figure environment, an image-inclusion command or drawing code, so no image is delivered and the package declares no asset dependency.  Editorial formatting only, in addition: an AMS \texttt{amsart} preamble that restates the article's own declarations and macros used by the retained lines, namely the environments \texttt{theorem} on separate counters, together with the standard packages those lines need.  The retained environments print in this document as Theorem 1 in order of appearance, whereas the article prints the same items with the numbering of its own full document.  The complete native source of the article as distributed in the arXiv e-print for this exact version is bundled unchanged as \texttt{sources/195818/source0.tex}.
\begin{theorem}[Characterization of Invisible Dynamics]\label{thm:invisible}
For $X \in \RR_*^{n \times d}$ (i.e., $\rank(X) = d$), a vector field $f$ produces invisible dynamics ($\dot{P} = 0$) if and only if $f(X) = XA$ for some skew-symmetric matrix $A \in \so(d)$.
\end{theorem}

\begin{proof}
($\Leftarrow$) If $f(X) = XA$ with $A^\top = -A$, then:
\begin{equation}
\dot{P} = f(X) X^\top + X f(X)^\top = XAX^\top + XA^\top X^\top = XAX^\top - XAX^\top = 0
\end{equation}

($\Rightarrow$) Suppose $\dot{P} = f(X) X^\top + X f(X)^\top = 0$.
Let $G = X^\top X$, which is invertible since $\rank(X)=d$.
Left-multiply the identity by $(I - XG^{-1}X^\top)$ to project onto $\col(X)^\perp$:
\begin{equation}
0 = (I - XG^{-1}X^\top) f(X) X^\top
\end{equation}
Since $X^\top$ has full row rank, this implies
\begin{equation}
(I - XG^{-1}X^\top) f(X) = 0
\end{equation}
hence $f(X) \in \col(X)$ and there exists $B \in \RR^{d \times d}$ such that $f(X) = XB$.

Substituting back into $\dot{P}=0$ gives
\begin{equation}
0 = X(B + B^\top) X^\top
\end{equation}
and since $X$ has full column rank, this forces $B + B^\top = 0$, i.e.\ $B \in \so(d)$.
Taking $A = B$ completes the proof.
\end{proof}

\end{document}
